% =====================================================================
%  DROP-IN SECTIONS: medspaCy mechanism + Figs 23/24/25 analysis
%  Written against the real evaluation logs:
%    Logs/evaluate_MedspacyUmlsChecker_full57_20260818_222322.json
%    Logs/evaluate_EmbedKdeChecker_full57_20260819_232035.json
%    Logs/evaluate_AlignScoreChecker_20260818_203942.json
%  (same JSONs that data_jsons/fig23..fig25 and graph_maker.py render from)
%
%  Figure PDFs have already been copied to:
%    figures/NLP/fig25_medspacy_tp_fn_by_corruption_type.pdf
%    figures/kde/fig23_kde_tp_fn_by_corruption_type.pdf
%    figures/align_score/fig24_alignscore_tp_fn_by_corruption_type.pdf
%
%  BLOCK A -> replaces main.tex line 760-762 (the one-paragraph mechanism)
%  BLOCK B -> replaces main.tex line 763 (the one-line results sentence)
%  BLOCK C -> insert into 9.3.4 Results, after the Table 4 paragraph
%  BLOCK D -> insert into 9.4.4 Results, after the Table 5 paragraph
%  BLOCK E -> optional cross-method table, for 9.5 / the Results chapter
% =====================================================================


% =====================================================================
% BLOCK A  --  how the checker works (mechanism, bullets)
% =====================================================================

The first SOAP checker module built on top of medspaCy reduces both documents to
the same machine readable object and then compares those two objects rather than
the text itself. The pipeline runs in three stages, and what each stage throws
away matters as much as what it keeps.

\paragraph{Stage 1 --- concept extraction.}
The shared QuickUMLS matcher \cite{soldaini2016quickumls} is run over the
transcript and over the SOAP note independently, and every span it resolves to a
UMLS Concept Unique Identifier (CUI) is kept, provided the concept falls inside
the semantic types this project scoped the matcher to:
\begin{itemize}
    \item \textbf{Diseases and findings} --- T047 (disease or syndrome), T033
          (finding), T184 (sign or symptom), T191 (neoplastic process).
    \item \textbf{Procedures} --- T059 (laboratory procedure), T060 (diagnostic
          procedure), T061 (therapeutic or preventive procedure).
    \item \textbf{Drugs} --- T121 (pharmacologic substance), T195 (antibiotic),
          T200 (clinical drug).
\end{itemize}
Any span outside those types, and any span the matcher cannot resolve at all, is
simply not carried forward; it does not exist for the rest of the pipeline.

\paragraph{Stage 2 --- assertion tagging.}
medspaCy's ConText component is then run over every surviving span and stamps it
with its assertion state, so that a mentioned concept and an \emph{asserted}
concept are not confused with one another:
\begin{itemize}
    \item \textbf{negated} --- ``no chest pain''
    \item \textbf{uncertain} --- ``possibly a chest infection''
    \item \textbf{historical} --- ``she had pneumonia in 2019''
    \item \textbf{hypothetical} --- ``if the rash spreads''
    \item \textbf{family associated} --- ``her mother has eczema''
\end{itemize}

\paragraph{Stage 3 --- set comparison.}
After the first two stages each document has been collapsed into a
\emph{set of (CUI, assertion) pairs}. The checker never looks at the prose
again; it compares the two sets by CUI and raises one of three deterministic
judgments, shown in \Cref{tab:medspacy_decision}.

\begin{table}[htbp]
  \centering
  \begin{tabular}{lll}
    \toprule
    \textbf{CUI is in transcript?} & \textbf{CUI is in note?} & \textbf{Judgment} \\
    \midrule
    yes & no  & \textbf{Omission} --- content the clinician said was dropped \\
    no  & yes & \textbf{Hallucination} --- content the note invented \\
    yes & yes & assertion states agree $\rightarrow$ no flag \\
    yes & yes & no pair of mentions agrees $\rightarrow$ \textbf{Status flip} \\
    \bottomrule
  \end{tabular}
  \caption{The three deterministic judgments of the medspaCy/UMLS/ConText checker.
  Every decision is a set membership test on CUIs; the status flip case is the only
  one that consults ConText's assertion state at all.}
  \label{tab:medspacy_decision}
\end{table}

\paragraph{What the comparison can and cannot see.}
Because the comparison happens on the reduced set and not the text, the checker
is structurally blind to anything that is not a CUI:
\begin{itemize}
    \item \textbf{Visible:} which clinical concepts appear on each side, and
          whether each mention of them is asserted, denied, hedged, historical,
          hypothetical or attributed to a relative.
    \item \textbf{Invisible:} \emph{numbers of every kind} --- doses, frequencies,
          durations, vital signs --- because a quantity is not a UMLS concept and
          carries no CUI. ``500\,mg three times a day'' and ``50\,mg once a day''
          reduce to the identical set.
    \item \textbf{Invisible:} word order, sentence boundaries, and which speaker
          said what. A doctor's question and a patient's one word answer land in
          the same undifferentiated bag as everything else.
    \item \textbf{Fragile:} near synonymous phrasing. A patient's ``loose stool''
          and a note's ``watery stools'' can resolve to two different CUIs under
          approximate matching, so a correctly documented fact still reads as a
          mismatch by concept identity alone.
    \item \textbf{Fragile:} ordinary conversational words that are themselves
          valid, if clinically generic, Metathesaurus entries --- ``well,''
          ``start,'' ``right'' --- which enter the transcript set as real
          concepts and then have nothing to match against in the note.
\end{itemize}
The last three bullets are the mechanism behind almost every number in
\Cref{fig:medspacy_by_type}; the second bullet is the mechanism behind one row of
it exactly.


% =====================================================================
% BLOCK B  --  9.2.2 results, Table 5 + Figure 25
% =====================================================================

\subsubsection{Results}
\label{sec:medspacy_results}

\begin{figure}[htbp]
    \centering
    \includegraphics[width=\textwidth]{figures/NLP/fig25_medspacy_tp_fn_by_corruption_type.pdf}
    \caption{medspaCy/UMLS/ConText checker: true positives and false negatives
    per injected corruption type over the full 57-file dataset, using the same
    six categories as \Cref{fig:label-distributions}. False positives carry no
    corruption type --- a flag that matches no ground-truth label has no label
    metadata to inherit --- so they are reported as a single untyped total on
    their own axis.}
    \label{fig:medspacy_by_type}
\end{figure}

Scored over the full prim 57-file evaluation dataset, this checker performs
poorly: 5.5\% precision and 16.7\% recall in the hallucination direction
(F1 = 8.3\%) \Cref{tab:results}, with omission detection catching nothing at all
(0\% precision and 0\% recall). Taken across all three directions at once the
figures are worse still --- 23 true positives against 2{,}193 false positives and
147 false negatives, or 1.0\% precision at 13.5\% recall --- and it is that
combined scope, not the hallucination only row of \Cref{tab:results}, that
\Cref{fig:medspacy_by_type} plots.

\Cref{fig:medspacy_by_type} breaks those 170 ground-truth labels down by the kind
of corruption that was injected, and the value of the breakdown is that it
separates a checker that is uniformly mediocre from one that is blind in
specific, diagnosable places. This checker is the second kind.

\begin{itemize}
    \item \textbf{Number edits: 0 of 33 caught.} Not a low score, a structurally
          impossible one. A dose, a frequency or a duration is not a UMLS concept
          and has no CUI, so a corrupted number and its original reduce to exactly
          the same $(\text{CUI},\,\text{assertion})$ pair and the set difference in
          \Cref{tab:medspacy_decision} is empty by construction. No threshold and
          no amount of tuning moves this row; only a different representation does.
    \item \textbf{Omitted detail: 0 of 32 caught, alongside 1{,}707 false positives
          in the same direction.} The omission judgment is not silent --- it is the
          single loudest part of the checker --- it is simply never right. Every
          conversational turn that resolves to a generic Metathesaurus entry becomes
          a transcript CUI with no counterpart in a short, condensed note, and is
          duly reported as lost clinical content. The direction fires constantly
          and scores nothing.
    \item \textbf{Entity swap: 1 of 23 caught (4\%).} The weakest non-zero row, and
          the one that most directly indicts CUI identity as a comparison key.
          Swapping one clinical entity for a neighbouring one usually lands inside
          the same approximate matching neighbourhood, so the swapped note concept
          still matches a transcript concept and no difference is raised.
    \item \textbf{Negation flip: 9 of 38 caught (24\%).} This is the most
          instructive row in the figure, because negation is the one failure mode
          medspaCy is specifically built for and gives a dedicated judgment to.
          Its dedicated status flip judgment fired 90 times across the 57 files and
          was correct \emph{zero} times. The 9 caught flips were all caught as
          \emph{hallucinations} instead --- the flip happened to change the CUI
          rather than only the assertion state, so the set difference caught what
          ConText did not. The component that was supposed to make this checker
          clinically aware contributed 90 false alarms and no correct flags.
    \item \textbf{Inserted sentence: 12 of 42 caught (29\%).} The best typed row,
          and the only category that plays to the design: a wholly fabricated
          sentence tends to introduce a genuinely new concept, which is precisely
          what a set difference is good at. Even here two thirds are missed, because
          a fabricated sentence assembled from concepts already present elsewhere in
          the transcript produces no new CUI at all.
    \item \textbf{Drug-name switch: 1 of 2 caught.} $n = 2$; nothing can be
          concluded from this row and it should not be read as a 50\% result.
\end{itemize}

The right hand panel of \Cref{fig:medspacy_by_type} is on its own axis for a
reason: at 2{,}193 false positives against 23 true positives, plotting it on the
shared scale would flatten every typed bar to invisibility. That ratio ---
roughly \textbf{95 false alarms for every correct flag} --- is the number that
disqualifies the module clinically, not the F1. A clinician working through this
checker's output would dismiss ninety-odd wrong flags before meeting a real one,
which is the alert fatigue failure mode described in \Cref{sec:Results} arriving
in its most extreme form.


% =====================================================================
% BLOCK C  --  9.3.4 KDE, Figure 23 + comparison to medspaCy
% =====================================================================

\begin{figure}[htbp]
    \centering
    \includegraphics[width=\textwidth]{figures/kde/fig23_kde_tp_fn_by_corruption_type.pdf}
    \caption{EmbedKDECheck: true positives and false negatives per injected
    corruption type over the same 57 files and the same 170 ground-truth labels
    as \Cref{fig:medspacy_by_type}, so the two figures can be read directly
    against each other.}
    \label{fig:kde_by_type}
\end{figure}

\Cref{fig:kde_by_type} plots the same breakdown for EmbedKDECheck, and read
against \Cref{fig:medspacy_by_type} it shows that dropping the symbolic layer
entirely bought a real improvement in every direction the method is capable of
looking:

\begin{itemize}
    \item \textbf{Every typed row improves, several of them severalfold.}
          Inserted sentences rise from 12 to 24 of 42, negation flips from 9 to 19
          of 38, and entity swaps from 1 to 14 of 23 --- a fourteenfold increase on
          the category where CUI identity failed hardest. Total typed catches rise
          from 23 to 76, lifting overall recall from 13.5\% to 44.7\%
          \Cref{tab:results}.
    \item \textbf{Number edits go from 0 to 17 of 33.} This is the clearest
          demonstration that the two methods fail for unrelated reasons. medspaCy's
          zero was structural: numbers have no CUI, so they cannot be compared.
          EmbedKDECheck has no such blind spot because it never reduces text to
          concepts at all --- it scores a whole sentence's tokens against a density
          map. The caveat is that it is not reading the number either; it is
          registering that the sentence containing it sits in a low density region,
          which is a weaker claim than ``this dose is wrong'' and explains why the
          row is a coin flip rather than a solved category.
    \item \textbf{Entity swaps are where being phrasing-agnostic pays.} A swapped
          entity that stayed inside the same UMLS neighbourhood, and so was
          invisible to a CUI set difference, still moves in embedding space far
          enough to change the density estimate. This is the single row that best
          justifies the move from rule-based NLP to a statistical method.
    \item \textbf{Omitted detail is still 0 of 32.} The one row that does not move.
          Both checkers score exactly nothing here, and --- damningly --- both do it
          while flagging omissions constantly: 1{,}707 omission-direction false
          positives for medspaCy, 1{,}330 for EmbedKDECheck. Two methods with
          nothing in common mechanically produce the same result in this direction,
          which is strong evidence that the problem is neither symbolic matching nor
          density estimation but the task itself. Scoring ordinary spoken
          conversational turns against a condensed note guarantees most of them look
          unsupported whether or not anything was actually lost, exactly the failure
          shown in \Cref{tab:KDE_sub_results}.
    \item \textbf{The false positive panel barely improves.} 1{,}661 against
          medspaCy's 2{,}193 --- a 24\% reduction while true positives more than
          tripled, which moves the cost per correct flag from roughly 95:1 to
          roughly 22:1. Better by a factor of four, and still four times too noisy
          to put in front of a clinician.
\end{itemize}

The honest summary of the medspaCy $\rightarrow$ KDE step is therefore that it
solved a recall problem and left a precision problem untouched. Removing the
requirement that a fact be phrased in a way a rule can recognise roughly tripled
what the checker could find; it did nothing about the fact that both methods are
being asked to compare a spoken register against a written one, and the false
positive panel is where that unsolved half of the problem shows up.


% =====================================================================
% BLOCK D  --  9.4.4 AlignScore, Figure 24 + comparison to both
% =====================================================================

\begin{figure}[htbp]
    \centering
    \includegraphics[width=\textwidth]{figures/align_score/fig24_alignscore_tp_fn_by_corruption_type.pdf}
    \caption{AlignScore (fixed pipeline, full 57-file run): true positives and
    false negatives per injected corruption type, over the same 170 ground-truth
    labels as \Cref{fig:medspacy_by_type} and \Cref{fig:kde_by_type}. Note that
    the false positive axis is an order of magnitude smaller than in either
    earlier figure.}
    \label{fig:alignscore_by_type}
\end{figure}

\Cref{fig:alignscore_by_type} completes the set, and the first thing to note is
what did \emph{not} change. AlignScore finds 70 of the 170 labels against
EmbedKDECheck's 76: on raw recall the trained NLI model is very slightly
\emph{worse} than the statistical density method it replaced. The entire gain
recorded in \Cref{tab:results} --- 8\% F1 to 30\% F1 --- is a precision gain, and
it is visible only in the right hand panel.

\paragraph{Where it is strongest.}
\begin{itemize}
    \item \textbf{False positives: 226, against EmbedKDECheck's 1{,}661 and
          medspaCy's 2{,}193.} Roughly \textbf{3 false alarms per correct flag},
          down from 22:1 and 95:1. Across the three-method progression the cost per
          catch falls by a factor of thirty while recall stays flat, which is the
          single clearest result the three figures produce together. Part of this
          is architectural rather than earned: because only the hallucination
          direction is implemented, AlignScore never generates the flood of
          omission-direction false positives that dominates both earlier panels.
          Even setting the omission direction aside entirely, though, it is still
          the only one of the three whose precision is in a plausible range at all.
    \item \textbf{Number edits: 20 of 33, the best row of any method.} This is the
          one place where the alignment head is genuinely doing something neither
          other method can. medspaCy scored zero here because quantities have no
          CUI; EmbedKDECheck scored 17 by accident of sentence-level density.
          AlignScore is the only one of the three that actually reads a quantity as
          part of the claim and can find a transcript span that fails to support it.
    \item \textbf{Entity swap (13 of 23) and inserted sentence (21 of 42)} hold
          roughly level with EmbedKDECheck, at a fraction of the false positive
          cost.
\end{itemize}

\paragraph{Where it is weakest.}
\begin{itemize}
    \item \textbf{Negation flips: 15 of 38, worse than EmbedKDECheck's 19.} The
          worst result in the figure relative to expectation, since a 3-way NLI
          head is trained precisely to separate entailment from contradiction and
          a negation flip is a textbook contradiction. The cause is the one
          diagnosed earlier in this section: the transcript-side chunking glues
          conditionally opposite instructions into a single $\sim$350-token context,
          so a chunk frequently supports both a claim and its negation and the
          alignment head has no reason to score either below threshold. The model is
          capable of this category; the surrounding pipeline denies it the evidence.
    \item \textbf{Omitted detail: 0 of 32, for the third time.} Here the zero is
          neither a representational blind spot (medspaCy) nor a register mismatch
          (EmbedKDECheck) but a straightforward absence: the symmetric direction is
          not implemented, so all 32 labels are automatic false negatives. This is
          also why AlignScore's omission direction contributes 0 false positives
          where the other two contributed 1{,}707 and 1{,}330 --- a meaningful part
          of its precision advantage is that it declines to answer the question the
          other two answer badly.
    \item \textbf{Inserted sentences still only half caught (21 of 42),} on the
          category that should be easiest: a fabricated claim with no transcript
          support is the exact profile AlignScore was trained for.
\end{itemize}

Read as a series, \Cref{fig:medspacy_by_type,fig:kde_by_type,fig:alignscore_by_type}
tell one story in three panels. Recall moved once, at the symbolic to statistical
step, and then stopped: 23 $\rightarrow$ 76 $\rightarrow$ 70 true positives. Precision
moved once, at the statistical to trained-NLI step, and moved by an order of
magnitude: 2{,}193 $\rightarrow$ 1{,}661 $\rightarrow$ 226 false positives. And one row
--- omitted detail --- is zero in all three, for three entirely different and
independently sufficient reasons. Three methods sharing no mechanism failing
identically on the same category is not three coincidences; it is the register
mismatch argument of \Cref{sec:Results} in its most direct empirical form.


% =====================================================================
% BLOCK E  --  optional cross-method summary table
% =====================================================================

\begin{table}[htbp]
  \centering
  \begin{tabular}{lrrrr}
    \toprule
    \textbf{Corruption type} & \textbf{Labels} & \textbf{medspaCy} & \textbf{EmbedKDECheck} & \textbf{AlignScore} \\
    \midrule
    Inserted sentence (hallucinated) & 42  & 12 (29\%) & \textbf{24 (57\%)} & 21 (50\%) \\
    Negation flip                    & 38  & 9 (24\%)  & \textbf{19 (50\%)} & 15 (39\%) \\
    Number edit                      & 33  & 0 (0\%)   & 17 (52\%)          & \textbf{20 (61\%)} \\
    Omitted detail                   & 32  & 0 (0\%)   & 0 (0\%)            & 0 (0\%) \\
    Entity swap                      & 23  & 1 (4\%)   & \textbf{14 (61\%)} & 13 (57\%) \\
    Drug-name switch                 & 2   & 1         & \textbf{2}         & 1 \\
    \midrule
    \textbf{True positives}          & 170 & 23        & \textbf{76}        & 70 \\
    \textbf{False positives}         & --- & 2{,}193   & 1{,}661            & \textbf{226} \\
    \textbf{False alarms per catch}  & --- & 95:1      & 22:1               & \textbf{3.2:1} \\
    \bottomrule
  \end{tabular}
  \caption{Recall per injected corruption type across the three cross-document
  checkers, all scored on the same 57 files and the same 170 hand-labelled
  ground-truth items (\Cref{fig:medspacy_by_type,fig:kde_by_type,fig:alignscore_by_type}).
  Drug-name switch is reported as raw counts; $n = 2$ is too small for a percentage.}
  \label{tab:checker_by_corruption_type}
\end{table}
